\documentclass{ppgeesa}

%\usepackage[latin1]{inputenc}
\usepackage[utf8]{inputenc}
\usepackage{graphicx,amsmath,amssymb,marvosym,latexsym,lineno,hyperref}
\usepackage{float}
\usepackage{arydshln}
\usepackage{bbm}
\usepackage{dsfont}
\graphicspath{{images}}
\usepackage{adjustbox}

\usepackage{tikz}
\usepackage{tikz-3dplot}
\usetikzlibrary{arrows.meta}
\usepackage{float}
\usepackage{adjustbox}
\usepackage{breqn}

\newtheorem{lema}{Lema}
\newcommand{\R}{\mathds{R}}

\begin{document}

\title{Modelagem e Linearização de Guindastes do tipo Ponte}

\author{
  Joe Ferreira Scholtz\\
	{Universidade Federal do Rio Grande do Sul (UFRGS) - Av. Osvaldo Aranha, 103 - Porto Alegre - Brasil} \\
}

\maketitle
\thispagestyle{empty}\pagestyle{empty}

\begin{abstract} % O abstract não deve exceder 100 palavras
Guindastes do tipo ponte rolante ou pórtico são sistemas amplamente utilizados em diversoes setores industriais para o transporte de cargas pesadas dentro de um espaço de trabalho e consistem em um sistema mecânico com três eixos de atuação e seis graus de liberdade, ou seja, um sistema subatuado. O Presente trabalho aborda a modelagem e linearização em torno de ponto de equilíbrio desse tipo de sistema, as condições de funcionamento e divergência em relação ao modelo não linear.
\end{abstract} % O abstract não deve exceder 100 palavras

\begin{IEEEkeywords}
Ponte Rolante, Pórtico, Modelagem de sistemas mecânicos, Linearização em torno de ponto de equilíbrio
\end{IEEEkeywords}

\section{Introdução}
Técnicas de visão computacional vem sendo amplamente utilizadas para desenvolver sistemas autônomos utilizando VANTs \cite{meyer2000matrix}.

\section{Modelagem do sistema}

Uma ponte rolante, representada na Figura \ref{fig:crane_tikz} , pode ser vista como o seguinte conjunto de estruturas: 

\begin{itemize}
  \item Trilho: Atuador linear ligado a uma estrutura fixa que move a Ponte em uma dada direção, nesse caso, a do eixo $X$;
  \item Ponte: Conectada à movimentação do trilho, é um atuador linear que move o Carro em uma direção perpedincular à direção do Trilho, nesse caso, a do eixo $Y$;
  \item Carro: Movimentado pela Ponte, é o ponto de origem do Cabo de Içamento;
  \item Cabo de Içamento: Cabo retrátil de comprimento variável $l$ que conecta o Carro à Carga;
  \item Carga: Corpo rígido de massa $m$ conectada à extremidade livre do Cabo de Içamento.
\end{itemize}

Sejam os sistemas de coordenadas, onde $(x,y,z)_{S}$ denota um ponto no espaço representado no sistema de coordenadas $S$.

\begin{itemize}
  \item $S_0 : \{X,Y,Z\}$: Sistema de coordenadas inercial global, fixo à estrutura do trilho, em uma de suas extremidades. Assume-se este sistema de coordenadas quando nenhum é explicitado;
  \item $S_C : {X_C,Y_C,Z_C}$: Sistema de coordenadas local alinhado à posição do Carro, com origem $(x,y,0)_{S0}$ e mesma orientação do sistema $S_0$.
\end{itemize}

\begin{figure}[H]
\begin{adjustbox}{width=\linewidth}
    \tdplotsetmaincoords{75}{115}
  \begin{tikzpicture}[tdplot_main_coords, line join=round, line cap=round]

      % --- Custom Macro to draw a 3D Box ---
      % Usage: \drawbox{x}{y}{z}{width_X}{width_Y}{height_Z}{base_color}
      \newcommand{\drawbox}[7]{%
        \def\X{#1} \def\Y{#2} \def\Z{#3}
        \def\dX{#4} \def\dY{#5} \def\dZ{#6}
        \colorlet{boxcol}{#7} 
        
        % BACK FACES (Drawn first so they are hidden by front faces)
        % Bottom face
        \draw[fill=boxcol!30!black, draw=black, semithick] 
          (\X, \Y, \Z) -- (\X+\dX, \Y, \Z) -- (\X+\dX, \Y+\dY, \Z) -- (\X, \Y+\dY, \Z) -- cycle;
        % -X face
        \draw[fill=boxcol!40!black, draw=black, semithick] 
          (\X, \Y, \Z) -- (\X, \Y+\dY, \Z) -- (\X, \Y+\dY, \Z+\dZ) -- (\X, \Y, \Z+\dZ) -- cycle;
        % -Y face
        \draw[fill=boxcol!50!black, draw=black, semithick] 
          (\X, \Y, \Z) -- (\X+\dX, \Y, \Z) -- (\X+\dX, \Y, \Z+\dZ) -- (\X, \Y, \Z+\dZ) -- cycle;
          
        % FRONT FACES (Drawn last to cover the back faces)
        % +X face (Faces front-left in this projection)
        \draw[fill=boxcol!70!white, draw=black, semithick] 
          (\X+\dX, \Y, \Z) -- (\X+\dX, \Y+\dY, \Z) -- (\X+\dX, \Y+\dY, \Z+\dZ) -- (\X+\dX, \Y, \Z+\dZ) -- cycle;
        % +Y face (Faces front-right in this projection)
        \draw[fill=boxcol, draw=black, semithick] 
          (\X, \Y+\dY, \Z) -- (\X+\dX, \Y+\dY, \Z) -- (\X+\dX, \Y+\dY, \Z+\dZ) -- (\X, \Y+\dY, \Z+\dZ) -- cycle;
        % +Z face (Faces top)
        \draw[fill=boxcol!40!white, draw=black, semithick] 
          (\X, \Y, \Z+\dZ) -- (\X+\dX, \Y, \Z+\dZ) -- (\X+\dX, \Y+\dY, \Z+\dZ) -- (\X, \Y+\dY, \Z+\dZ) -- cycle;
      }

      % --- Variables for robot position ---
      \def\MaxXpos{5} % Max x value
      \def\Xpos{2.5}  % Position along X-axis
      \def\MaxYpos{5} % Max x value
      \def\Ypos{2.5}  % Position along Y-axis
      \def\MaxZpos{5} % Max x value
      \def\Zpos{1.2}  % Position along Z-axis

            % --- Pendulum Parameters ---
      \def\PendL{4.0}  % Length of the pendulum
      \def\ThetaX{28}  % Swing angle in XZ plane (degrees)
      \def\ThetaY{35}  % Swing angle in YZ plane (degrees)
      \def\ArcR{2.0}   % Radius for the drawn angle arcs

      % compute mass coordinates
      % tan(theta_x) = X/(-Z) and tan(theta_y) = Y/(-Z)
      \pgfmathsetmacro{\TanX}{tan(\ThetaX)}
      \pgfmathsetmacro{\TanY}{tan(\ThetaY)}
      \pgfmathsetmacro{\Denom}{sqrt(\TanX*\TanX + \TanY*\TanY + 1)}
      \pgfmathsetmacro{\Dz}{-\PendL / \Denom}
      \pgfmathsetmacro{\Dx}{-\Dz * \TanX}
      \pgfmathsetmacro{\Dy}{-\Dz * \TanY}
      
      % Final absolute coordinates for the pendulum mass
      \pgfmathsetmacro{\Bx}{\Xpos + \Dx}
      \pgfmathsetmacro{\By}{\Ypos + \Dy}
      \pgfmathsetmacro{\Bz}{\Dz}



      % --- Global Reference ---
      \def\ArrowType{{-Stealth[scale=1.5]}}
      \draw[dashed,black] (0,0,0) -- +(-4,0,0);
      \draw[dashed,black] (0,0,0) -- +(0,-3,0);
      \draw[dashed,black] (0,0,0) -- +(0,0,-2);
      % measures
      \draw[dashed,black] (\Xpos,0,0) -- +(0,-3,0);
      \draw[dashed,Stealth-Stealth,black] (0,-1.5,0) -- +(\Xpos,0,0) node[midway,above] {$x$};
      \draw[dashed,black] (\Xpos,\Ypos,0) -- +(-\Xpos-4,0,0);
      \draw[dashed,Stealth-Stealth,black] (-2,0,0) -- +(0,\Ypos,0) node[midway,above] {$y$};


      % --- Trilho ---
      \def\TrilhoDx{0.5}
      \def\TrilhoLy{0.4}
      \def\TrilhoLz{0.15}
      % Bounding Box
      \draw[dashed, black, thin] (\Bx, \Ypos, 0) -- (\Bx, \Ypos, \Bz);
      \draw[dashed, black, thin] (\Bx, \By, -\TrilhoLz) -- (\Bx, \By, \Bz);
      \draw[dashed, black, thin] (\Bx, \By, -\TrilhoLz) -- (\Bx, \By, \Bz);
      \draw[dashed, black, thin] (\Xpos, \By, -\TrilhoLz) -- (\Xpos, \By, \Bz);
      \draw[dashed, black, thin] (\Bx, \Ypos, \Bz) -- (\Bx, \By, \Bz) -- (\Xpos, \By, \Bz);
      \draw[dashed, black, thin] (\Xpos, \Ypos, \Bz) -- (\Bx, \Ypos, \Bz);
      \draw[dashed, black, thin] (\Xpos, \Ypos, \Bz) -- (\Xpos, \By, \Bz);
      %\drawbox{x}{y}{z}{width_X}{width_Y}{height_Z}{base_color}
      \drawbox{-\TrilhoDx}{-\TrilhoLy*0.5}{-\TrilhoLz}{\MaxXpos+2*\TrilhoDx}{\TrilhoLy}{\TrilhoLz}{black!60}
      \drawbox{-\TrilhoDx}{\MaxYpos-\TrilhoLy*0.5}{-\TrilhoLz}{\MaxXpos+2*\TrilhoDx}{\TrilhoLy}{\TrilhoLz}{black!60}

      % --- Global Reference ---
      \def\ArrowType{{-Stealth[scale=1.5]}}
      \draw[thick,\ArrowType, black] (0,0,0) -- (8,0,0) node[anchor=north east]{$X$};
      \draw[thick,\ArrowType, black] (0,0,0) -- (0,6,0) node[anchor=north west]{$Y$};
      \draw[thick,\ArrowType, black] (0,0,0) -- (0,0,2) node[anchor=south]{$Z$};

      % --- Ponte ---
      %\drawbox{x}{y}{z}{width_X}{width_Y}{height_Z}{base_color}
      \def\PonteLx{0.8}
      \def\PonteDy{0.2}
      \def\PonteLz{0.2}
      % \drawbox{\Xpos-0.5*\PonteLx}{-\PonteDy}{0}{\PonteLx}{\MaxYpos+2*\PonteDy}{\PonteLz}{black!60}
      \drawbox{\Xpos-0.5*\PonteLx}{-\PonteDy}{0}{\PonteLx}{\Ypos+\PonteDy}{\PonteLz}{black!60}


      % 1. Vertical reference line (Zero angle)
      \draw[dashed, black, thin] (\Xpos, \Ypos, 0) -- (\Xpos, \Ypos, \Bz);

      % 2. XZ Plane Projection
      % Arc for \theta_x (parameterized plot acting as an arc in 3D)
      \begin{scope}[shift={(\Xpos,\Ypos,0)}, x={(1,0,0)}, y={(0,0,1)}]
        \draw[black, thin, -Stealth] (0,-\ArcR) arc (-90:-90+\ThetaX:\ArcR);
        \def\LabelTxFraction{0.5}
        \node[black] at ({-90+\ThetaX*\LabelTxFraction}:\ArcR*1.15) {$\theta_x$};
      \end{scope}

      % 3. YZ Plane Projection (Red)
      \draw[dashed, black, thin] (\Xpos, \Ypos, 0) -- (\Xpos, \By, \Bz);
      \draw[dashed, black, thin] (\Xpos, \Ypos, \Bz) -- (\Bx, \By, \Bz);
      % Arc for \theta_y (Also removed 'smooth' here to fix its arrow angle)
      \draw[black, thin, -Stealth] plot[domain=0:\ThetaY, samples=40, variable=\t]
        ({\Xpos}, {\Ypos + \ArcR*sin(\t)}, {-\ArcR*cos(\t)});
      % Change \LabelFraction to adjust where the label sits along the arc (0 to 1)
      \def\LabelTyFraction{0.4}
      \node[black] at ({\Xpos}, {\Ypos + \ArcR*1.15*sin(\ThetaY*\LabelTyFraction)}, {-\ArcR*1.15*cos(\ThetaY*\LabelTyFraction)}) {$\theta_y$};

      % 5. The actual Pendulum Rod and Bob
      \draw[line width=0.5mm, black] (\Xpos, \Ypos, 0) -- (\Bx, \By, \Bz) node[midway,above,xshift=-8pt] {$l$};
      \draw[red!50!black,very thick,-Latex] (\Bx, \By, \Bz) -- +(0,0,-1) node[below] {$F_P = mg$};
      \node[circle, minimum size=0.5cm, inner sep=0pt, ball color=black!50!white, draw=black, semithick] 
        at (\Bx, \By, \Bz) {};

      % --- Carro ---
      %\drawbox{x}{y}{z}{width_X}{width_Y}{height_Z}{base_color}
      \def\CarroLx{1.2}
      \def\CarroLy{0.8}
      \def\CarroLz{0.4}
      \drawbox{\Xpos-0.5*\CarroLx}{\Ypos-0.5*\CarroLy}{0}{\CarroLx}{\CarroLy}{\CarroLz}{black!60}

      % --- Ponte ---
      %\drawbox{x}{y}{z}{width_X}{width_Y}{height_Z}{base_color}
      \drawbox{\Xpos-0.5*\PonteLx}{\Ypos+0.5*\CarroLy}{0}{\PonteLx}{\MaxYpos-\Ypos-\CarroLy*0.5+\PonteDy}{\PonteLz}{black!60}

      % --- Local Reference ---
      \def\ArrowType{{-Stealth[scale=1.5]}}
      \draw[dotted,black] (\Xpos,\Ypos,0) -- +(8-\Xpos,0,0);
      \draw[dotted,black] (\Xpos,\Ypos,0) -- +(0,6-\Ypos,0);
      \draw[dotted, black] (\Xpos,\Ypos,0) -- +(0,0,2);
      % --- Local Reference ---
      \def\ArrowType{{-Stealth[scale=1.5]}}
      \draw[thick,\ArrowType, black] (\Xpos+0.5*\CarroLx,\Ypos,0) -- +(8-\Xpos,0,0) node[anchor=north east]{$X_C$};
      \draw[thick,\ArrowType, black] (\Xpos,\MaxYpos+0.5*\TrilhoDx,0) -- +(0,6-\MaxYpos,0) node[anchor=north west]{$Y_C$};
      \draw[thick,\ArrowType, black] (\Xpos,\Ypos,\CarroLz) -- +(0,0,2) node[anchor=south]{$Z_C$};

      \draw[dashed, Stealth-Stealth,black, thin] (\Bx, \Ypos, \Bz) -- (\Xpos, \Ypos, 0) node[midway,above, left=-1pt] {};

      % --- Global Reference ---
      \def\ArrowType{{-Stealth[scale=1.5]}}
      \draw[dotted, black] (0,0,0) -- (8,0,0);
      \draw[dotted, black] (0,0,0) -- (0,6,0);
      \draw[dotted, black] (0,0,0) -- (0,0,2);



      % --- Component Labels ---
      % Using >=Stealth ensures the arrow head matches your other arrows
      % The bend left/right commands draw nice curved lines to the labels
      \draw[thin, <-, >=Stealth,black!80!white] (\MaxXpos+1, -0.1, 0.1) 
        to[bend left] ++(1.0, -0.5, 0.8) node[above] {Trilho};
        
      \draw[thin, <-, >=Stealth,black!80!white] (\Xpos, \MaxYpos-0.5, \PonteLz) 
        to[bend left] ++(-3.2, -0.8, 0.1) node[right] {Ponte};
        
      \draw[thin, <-, >=Stealth,black!80!white] (\Xpos-0.2*\CarroLx, \Ypos, \CarroLz) 
        to[bend left] ++(-4.8, -0.8, 0.5) node[right] {Carro};
        
      % Points just to the right of the ball's center so it touches the edge
      \draw[thin, <-, >=Stealth,black!80!white] (\Bx+0.8, \By+0.6, \Bz) 
        to[bend right] ++(1.5, 1.5, 0) node[right] {Carga};

      \draw[thin, <-, >=Stealth,black!80!white] (\Bx, \By-0.3, \Bz+1) 
        to[bend left] ++(1.5, 2.5, 1) node[right] {Cabo de Içamento};

  \end{tikzpicture}

\end{adjustbox}
\caption{Representação de uma ponte rolante tridimensional, suas estruturas e os sistemas de coordenadas adotados. Fonte: Adaptado de \cite{thiago2026}.}
\label{fig:crane_tikz}
\end{figure}


O vetor de configuração do sistema é dado como

\begin{equation}
\mathbf{q}
= 
\left[
\begin{array}{c}
x\\
y\\
l\\
\theta_x\\ 
\theta_y
\end{array} 
\right]\\
\end{equation}

Com $x$,$y$ posições do carro, e $l>0$ o comprimento do cabo de içamento da carga, as juntas atuadas, e $\theta_x$ e $\theta_y$ são os angulos da carga em relação, respectivamente, ao eixos $X$ e $Y$ de $S_0$, não atuados.

\subsection{Modelo Geométrico Direto}

Dessa forma, o modelo geométrico direto, isto é, a posição da carga em função das variáveis de junta, é dado por

\begin{align}
  \mathbf{p}_m(\mathbf{q})
&= 
\left[
\begin{array}{c}
x_m\\
y_m\\
z_m
\end{array} 
\right]\\
&=
\left[
\begin{array}{c}
x\\
y\\
0
\end{array} 
\right]
+
\left[
\begin{array}{c}
l\sin(\theta_x)\cos(\theta_y)\\
l\sin(\theta_y)\\
-l\cos(\theta_x)\cos(\theta_y)
\end{array} 
\right]
\label{eq:mass_position}
\end{align}

\subsection{Modelo Cinemático Direto}

O modelo cinemático direto, por sua vez, que relaciona a velocidade da carga com as derivadas das variáveis de junta no tempo fica dado por

\begin{equation}
\dot{\mathbf{p}}_m(\mathbf{q})
= 
\left[
  \begin{array}{c}
    \dot{x}_m\\
    \dot{y}_m\\
    \dot{z}_m
  \end{array} 
\right]
\end{equation}

com 

\begin{align}
\dot{x}_m &= \dot{x}+\dot{l}\sin(\theta_x)\cos(\theta_y)\nonumber\\
    &+ l\left(\cos(\theta_x)\cos(\theta_y)\dot{\theta}_x-\sin(\theta_x)\sin(\theta_y)\dot{\theta}_y\right)\\
\dot{y}_m &= \dot{y} + \dot{l}\sin(\theta_y) + l\cos(\theta_y)\dot{\theta}_y\\
\dot{z}_m &= -\dot{l}\cos(\theta_x)\cos(\theta_y)\nonumber\\
    &+ l\left(\sin(\theta_x)\cos(\theta_y)\dot{\theta}_x+\cos(\theta_x)\sin(\theta_y)\dot{\theta}_y\right)
\label{eq:direct_kinematics}
\end{align}

\subsection{Modelo Dinâmico}

Aqui são consideradas as seguintes simplificações:

\begin{itemize}
  \item A massa da carga é considerada pontual;
  \item A flexibilidade e massa do cabo são desprezadas;
  \item O atrito é considerado como puramente viscoso, diretamente proporcial à velocidade das juntas. 
\end{itemize}

\noindent e a abordagem adotada para derivação do modelo dinâmico segue Equação de Euler-Lagrange dada por

\begin{equation}
  \frac{d}{dt}\left(\frac{\partial L(\mathbf{q}, \dot{\mathbf{q}})}{\partial \dot{q_i}}\right) - \frac{\partial L(\mathbf{q}, \dot{\mathbf{q}})}{\partial q_i} = Q_{g,i}
\end{equation}

\noindent onde $L(\mathbf{q}, \dot{\mathbf{q}})$ é o Lagrangiano, definido por

\begin{equation}
  L(\mathbf{q}, \dot{\mathbf{q}}) \triangleq K(\mathbf{q}, \dot{\mathbf{q}}) - P(\mathbf{q})\label{eq:lagrangian}
\end{equation}

\noindent para $K(\mathbf{q}, \dot{\mathbf{q}})$ a energia cinética de todo o sistema e $P(\mathbf{q})$ a energia potencial total da carga, e $Q_{g,i}$ um esforço genérico aplicado à i-ésima junta.

Tal formulação, ainda que correta, não deixa explícita a influência de esforços não conservativos tais como atrito e opta por esforços genéricos $Q_{g,i}$ para os englobar. Visto que, na abordagem de modelagem adotada, é desejado considerar a influência dos atritos viscosos, opta-se por deixá-lo evidente. De fato, para sistemas cujos únicos esforços não conservativos são do tipo dissipativos proporcionais à velocidade, pode-se considerar $Q_{g,i}$ como tendo duas componentes como segue \cite{Meirovitch2012}: %pg 90

\begin{equation}
  Q_{g,i} = -\frac{\partial R(\dot{\mathbf{q}})}{\partial \dot{q_i}} + Q_{i}
\end{equation}

\noindent onde $Q_i$ é representa a componente do vetor de esforços externos relativa à i-ésima junta (e.g. esforço de atuação, ou perturbação na junta) e $R$ é dada pela função de dissipação de Rayleigh, caracterizada por

\begin{equation}
  R(\dot{\mathbf{q}}) = \frac{1}{2}\left(D_x\dot{x}^2+D_y\dot{y}^2+D_l\dot{l}^2\right)\label{eq:rayleigh}
\end{equation}

\noindent e que modela aqui o atrito viscoso do sistema. Onde $D_x$, $D_y$ e $D_l$ representam os coeficientes de amortecimento viscoso associados, respectivamente, ao movimentos nas direções $x,y,l$ e dependem unicamente no sistema de coordenadas e não do valor das velocidades.

Dessa forma, é possível reescrever a Equação de Euler-Lagrange para sistemas não conservativos, cujas forças não conservativas são do tipo dissipativas proporcionais à velocidade como \cite{thiago2026}:

\begin{equation}
  \frac{d}{dt}\left(\frac{\partial L(\mathbf{q}, \dot{\mathbf{q}})}{\partial \dot{q_i}}\right) - \frac{\partial L(\mathbf{q}, \dot{\mathbf{q}})}{\partial q_i} + \frac{\partial R(\dot{\mathbf{q}})}{\partial \dot{q_i}} = Q_{i}\label{eq:euler_lagrange_non_conservative}
\end{equation}

A fim de calcular o Lagrangiano, a energia cinética do sistema e potencial da carga podem serem reescritas como

\begin{align}
  K(\mathbf{q}, \dot{\mathbf{q}}) &= \frac{1}{2}(M_x \dot{x}^2 + M_y \dot{y}^2 + M_{l} \dot{l}^2) + \frac{m}{2} v_m^2 \label{eq:kinematic_en}\\
  P(\mathbf{q}) &= -m g l \cos(\theta_x) \cos(\theta_y)\label{eq:potential_en}
\end{align}

\noindent onde $M_x$ , $M_y$ e $M_l$ representam respectivamente as inércias equivalentes da ponte, carro e mecanismo de içamento
e $g$ é a aceleração da gravidade. Nesse caso, a velocidade da carga $v_m$ pode ser determinada como segue

\begin{align}
  v_m^2 &= \dot{x}_m^2 + \dot{y}_m^2 + \dot{z}_m^2 \nonumber \\ 
        & = \dot{x}^2 + \dot{y}^2 + \dot{l}^2 + l^2\cos^2(\theta_y)\dot{\theta}_x^2 + l^2\dot{\theta}_y^2 \nonumber\\
        &+2\left(\sin(\theta_x)\cos(\theta_y)\dot{l}+ l\cos(\theta_x)\cos(\theta_y)\dot{\theta}_x\right. \nonumber\\
        &- \left. l\sin(\theta_x)\sin(\theta_y)\dot{\theta}_y\right)\dot{x} \nonumber\\
        &+2\left(sin(\theta_y)\dot{l}+l\cos(\theta_y)\dot{\theta}_y\right)\dot{y}\label{eq:vm}
\end{align}

Adotando $Q_i = \left[\begin{array}{ccccc}F_x &F_y &F_l &0 &0\end{array}\right]^T$ -- ou seja, forças externas apenas provenientes dos atuadores de junta, sem considerar perturbações-- podemos finalmente, a partir das equações \ref{eq:lagrangian}, \ref{eq:rayleigh}, \ref{eq:euler_lagrange_non_conservative}, \ref{eq:kinematic_en}, \ref{eq:potential_en} e \ref{eq:vm} derivar o modelo dinâmico não linear desse sistema, que resulta nas seguintes expressões\cite{lee}

\begin{align}
F_x &=\nonumber\\
&\begin{pmatrix}
(M_x + m)\ddot{x} + ml\cos(\theta_x)\cos(\theta_y)\ddot{\theta}_x - ml\sin(\theta_x)\sin(\theta_y)\ddot{\theta}_y \\
+ m\sin(\theta_x)\cos(\theta_y)\ddot{l} + D_x\dot{x} + 2m\cos(\theta_x)\cos(\theta_y)\dot{l}\dot{\theta}_x \\
- 2m\sin(\theta_x)\sin(\theta_y)\dot{l}\dot{\theta}_y - ml\sin(\theta_x)\cos(\theta_y)\dot{\theta}_x^2 \\
- 2ml\cos(\theta_x)\sin(\theta_y)\dot{\theta}_x\dot{\theta}_y - ml\sin(\theta_x)\cos(\theta_y)\dot{\theta}_y^2
\end{pmatrix}\\
F_y &=\nonumber\\
&\begin{pmatrix}
(M_y + m)\ddot{y} + ml \cos(\theta_y)\ddot{\theta}_y + m\sin(\theta_y)\ddot{l} \\
{}+ D_y\dot{y} + 2m\cos(\theta_y)\dot{l}\dot{\theta}_y - ml\sin(\theta_y)\dot{\theta}_y^2
\end{pmatrix}\\
F_l &=\nonumber\\
&\begin{pmatrix}
(M_l + m)\ddot{l} + m\sin(\theta_x)\cos(\theta_y)\ddot{x} + m\sin(\theta_y)\ddot{y} + D_l\dot{l} \\
{}- ml\cos^2(\theta_y)\dot{\theta}_x^2 - ml\dot{\theta}_y^2 - mg\cos(\theta_x)\cos(\theta_y)
\end{pmatrix}\\
0 &=\nonumber\\
&\begin{pmatrix}
ml^2\cos^2(\theta_y)\ddot{\theta}_x + ml\cos(\theta_x)\cos(\theta_y)\ddot{x} + 2ml\cos^2(\theta_y)\dot{l}\dot{\theta}_x \\
{}- 2ml^2\sin(\theta_y)\cos(\theta_y)\dot{\theta}_x\dot{\theta}_y + mgl\sin(\theta_x)\cos(\theta_y)
\end{pmatrix}\\
0 &=\nonumber\\
&\begin{pmatrix}
ml^2\ddot{\theta}_y + ml\cos(\theta_y)\ddot{y} - ml\sin(\theta_x)\sin(\theta_y)\ddot{x} + 2ml\dot{l}\dot{\theta}_y \\
{}+ ml^2\cos(\theta_y)\sin(\theta_y)\dot{\theta}_x^2 + mgl\cos(\theta_x)\sin(\theta_y)
\end{pmatrix}
\end{align}


\begin{align*}
\mathbf{F} &= \begin{bmatrix} F_x & F_y & F_l & 0 & 0 \end{bmatrix}^T, \\
\mathbf{G}(\mathbf{q}) &= \begin{bmatrix} 0 & 0 & g_3 & g_4 & g_5 \end{bmatrix}^T, \\
\mathbf{D} &= \begin{bmatrix}
D_x & 0 & 0 & 0 & 0 \\
0 & D_y & 0 & 0 & 0 \\
0 & 0 & D_l & 0 & 0 \\
0 & 0 & 0 & 0 & 0 \\
0 & 0 & 0 & 0 & 0
\end{bmatrix}, \\
\mathbf{C}(\mathbf{q}, \dot{\mathbf{q}}) &= \begin{bmatrix}
0 & 0 & c_{13} & c_{14} & c_{15} \\
0 & 0 & c_{23} & 0 & c_{25} \\
0 & 0 & 0 & c_{34} & c_{35} \\
0 & 0 & c_{43} & c_{44} & c_{45} \\
0 & 0 & c_{53} & c_{54} & c_{55}
\end{bmatrix}, \\
\mathbf{M}(\mathbf{q}) &= \begin{bmatrix}
m_{11} & 0 & m_{13} & m_{14} & m_{15} \\
0 & m_{22} & m_{23} & 0 & m_{25} \\
m_{31} & m_{32} & m_{33} & 0 & 0 \\
m_{41} & 0 & 0 & m_{44} & 0 \\
m_{51} & m_{52} & 0 & 0 & m_{55}
\end{bmatrix}.
\end{align*}


\begin{align*}
m_{11} &= M_x + m, \\
m_{13} &= m_{31} = m \sin(\theta_x) \cos(\theta_y), \\
m_{14} &= m_{41} = ml \cos(\theta_x) \cos(\theta_y), \\
m_{15} &= m_{51} = -ml \sin(\theta_x) \sin(\theta_y), \\
m_{22} &= M_y + m, \\
m_{23} &= m_{32} = m \sin(\theta_y), \\
m_{25} &= m_{52} = ml \cos(\theta_y), \\
m_{33} &= M_l + m, \\
m_{44} &= ml^2 \cos^2(\theta_y), \\
m_{55} &= ml^2, \\
c_{13} &= m \cos(\theta_x) \cos(\theta_y) \dot{\theta}_x - m \sin(\theta_x) \sin(\theta_y) \dot{\theta}_y, \\
c_{14} &= m \cos(\theta_x) \cos(\theta_y) \dot{l} - ml \sin(\theta_x) \cos(\theta_y) \dot{\theta}_x \\
    &- ml \cos(\theta_x) \sin(\theta_y) \dot{\theta}_y, \\
c_{15} &= -m \sin(\theta_x) \sin(\theta_y) \dot{l} - ml \cos(\theta_x) \sin(\theta_y) \dot{\theta}_x \\
    &- ml \sin(\theta_x) \cos(\theta_y) \dot{\theta}_y, \\
c_{23} &= m \cos(\theta_y) \dot{\theta}_y, \\
c_{25} &= m \cos(\theta_y) \dot{l} - ml \sin(\theta_y) \dot{\theta}_y, \\
c_{34} &= -ml \cos^2(\theta_y) \dot{\theta}_x, \\
c_{35} &= -ml \dot{\theta}_y, \\
c_{43} &= ml \cos^2(\theta_y) \dot{\theta}_x, \\
c_{44} &= ml \cos^2(\theta_y) \dot{l} - ml^2 \sin(\theta_y) \cos(\theta_y) \dot{\theta}_y, \\
c_{45} &= -ml^2 \sin(\theta_y) \cos(\theta_y) \dot{\theta}_x, \\
c_{53} &= ml \dot{\theta}_y, \\
c_{54} &= ml^2 \cos(\theta_y) \sin(\theta_y) \dot{\theta}_x, \\
c_{55} &= ml \dot{l}, \\
g_3 &= -mg \cos(\theta_x) \cos(\theta_y), \\
g_4 &= mgl \sin(\theta_x) \cos(\theta_y), \\
g_5 &= mgl \cos(\theta_x) \sin(\theta_y).
\end{align*}

\subsection{Linearização do Modelo Dinâmico}

\begin{align*}
\mathbf{F}_{\text{lin}} &= \begin{bmatrix} F_x & F_y & F_l & 0 & 0 \end{bmatrix}^T, \\
\mathbf{G}_{\text{lin}}(\mathbf{q}) &= \begin{bmatrix} 0 & 0 & -mg & g\theta_x & g\theta_y \end{bmatrix}^T, \\
\mathbf{D}_{\text{lin}} &= \begin{bmatrix}
D_x & 0 & 0 & 0 & 0 \\
0 & D_y & 0 & 0 & 0 \\
0 & 0 & D_l & 0 & 0 \\
0 & 0 & 0 & 0 & 0 \\
0 & 0 & 0 & 0 & 0
\end{bmatrix}, \\
\mathbf{M}_{\text{lin}}(\mathbf{q}) &= \begin{bmatrix}
M_x + m & 0 & 0 & ml & 0 \\
0 & M_y + m & 0 & 0 & ml \\
0 & 0 & M_l + m & 0 & 0 \\
1 & 0 & 0 & l & 0 \\
0 & 1 & 0 & 0 & l
\end{bmatrix}.
\end{align*}


\subsection{Representação no Espaço de Estados do Modelo Dinâmico Linearizado}

\section{Resposta do Sistema em Malha Aberta}

Em trabalho futuro, esta seção irá abordar os seguintes tópicos:

\begin{enumerate}
  \item Resposta temporal à entrada zero do sistema em malha aberta para diferentes condições iniciais;
  \item Resposta temporal ao estado zero do sistema em malha aberta para referências de interesse;
  \item Apresentação dos diagramas de simulação em malha aberta e comparação dos modelos não linear e linearizado;
  \item Analise a estabilidade, observabilidade e controlabilidade do sistema; caso o sistema não seja controlável e/ou observável, será analisada a decomposição de Kalman do mesmo.
\end{enumerate}

\section{Projeto de Controlador e Observador de Estados}

Em trabalho futuro, esta seção irá abordar os seguintes tópicos:

\begin{enumerate}
  \item Formulação dos requisitos de controle;
  \item Determinação da abordagem de controle;
  \item Projeto de um controlador em malha fechada;
  \item Projeto de um observador de estados para implementação do controle proposto, justificando a escolha dos polos do observador em relação ao sistema em malha fechada;
  \item Apresentação dos diagramas de simulação em malha fechada;
  \item Análise dos resultados e comparação de desempenho em malha aberta e malha fechada sem e com observador de estados.
\end{enumerate}

\section{Conclusão}

Em trabalho futuro, esta seção discutirá de forma resumida tudo que foi abordado, desde os diferentes modelos do sistema, até o projeto de controle e os resultados obtidos.

\bibliographystyle{IEEEtran.bst}
\bibliography{base.bib}

\end{document}

